\begin{proof}
The two fixed points are distinct because $F'(r)<1$.
For a nonaffine map the second fixed point is real and outside $I$,
since the contraction has only one fixed point in $I$.
The map $\psi(t)=(t-r)/(t-\bar r)$ conjugates $F$ to multiplication
by $F'(r)$, giving \eqref{II:U:eq:rank-two-cross-ratio}.
Since $t_0,r$ lie on the same side of $\bar r$, $c<1$ and
$1-c\gamma^L\ge\min(1,1-c)>0$. One may use
$C_r=|r-\bar r||c|/\min(1,1-c)$, or $C_r=|t_0-r|$
in the affine case, and take the larger endpoint constant.
\end{proof}
